\section{Analytical Methods}
\label{sec:methods}

\subsection{Equations, Conditions, and Connection to the Application}

The three ODEs derived above--\eqref{eq:rc1st}, \eqref{eq:rl1st},
and~\eqref{eq:rlc2nd}--together model the spectrum from simple to complex
circuit behavior.
With initial conditions applied, the initial value problems (IVPs) are:

\noindent\textbf{RC circuit (first-order):}
\begin{equation}
  \frac{d\vC}{dt} + \frac{1}{RC}\,\vC = \frac{\Vs}{RC}, \qquad \vC(0) = V_0.
  \label{eq:rc_ivp}
\end{equation}

\noindent\textbf{RL circuit (first-order):}
\begin{equation}
  \frac{di}{dt} + \frac{R}{L}\,i = \frac{\Vs}{L}, \qquad i(0) = 0.
  \label{eq:rl_ivp}
\end{equation}

\noindent\textbf{RLC circuit (second-order):}
\begin{equation}
  LC\frac{d^2\vC}{dt^2} + RC\frac{d\vC}{dt} + \vC = \Vs,
  \qquad \vC(0) = \qty{0}{\volt},\quad \dot{\vC}(0) = 0.
  \label{eq:rlc_ivp}
\end{equation}

Each equation models a physical balance: the left-hand terms represent
energy dissipation (via $R$) and energy storage (via $L$ and $C$), while
the right-hand side represents the forcing from the voltage source.

\subsection{Independent and Dependent Variables}

\begin{itemize}
  \item \textbf{Independent variable:} time $t$ in \si{\second}.
  \item \textbf{Dependent variables:} $\vC(t)$ in \si{\volt} (capacitor
        voltage) and $i(t)$ in \si{\ampere} (series current).
\end{itemize}

\subsection{Parameters}

\begin{itemize}
  \item $R$: resistance; dissipates electrical energy as heat. Appears in the
        first-order damping term $R/L$ and the second-order damping
        coefficient $\alpha = R/(2L)$.
  \item $L = \qty{0.1}{\henry}$: inductance; stores energy in a magnetic field.
        Governs how quickly current can change.
  \item $C = \qty{10e-6}{\farad}$: capacitance; stores energy in an electric
        field. Governs how quickly voltage can change.
  \item $\Vs = \qty{10}{\volt}$: step-input amplitude. Represents a DC voltage
        switched on at $t = 0$.
\end{itemize}

\subsection{Assumptions and Limitations}

\paragraph{Assumptions.}
We adopt the standard \emph{lumped-element model}: all resistance is
concentrated in $R$, all inductance in $L$, and all capacitance in $C$,
with no parasitic coupling between components.
Component values are assumed to be constant with respect to temperature,
frequency, and signal amplitude (i.e., linear, time-invariant behavior).
Wires and switch contacts are treated as ideal (zero resistance, zero
inductance).

\paragraph{Limitations.}
Real resistors exhibit temperature-dependent resistance governed by
$R(T) = R_0[1 + \alpha_R(T - T_0)]$, introducing nonlinear behavior
under large current loads.
At high frequencies, parasitic inductance in capacitors and parasitic
capacitance in inductors violate the lumped-element assumption, requiring
distributed-element or transmission-line models.
The ideal step-input assumption ignores source impedance and switching
transients present in physical systems.
These effects are left for future work.

\subsection{Analytic Solutions}

\subsubsection{First-Order RC Circuit}

Equation~\eqref{eq:rc_ivp} is a linear first-order ODE with constant
coefficients.
Using the integrating factor $\mu(t) = e^{t/(RC)}$, the general solution is:
\begin{equation}
  \vC(t) = \Vs + (V_0 - \Vs)\,e^{-t/(RC)}.
  \label{eq:rc_sol}
\end{equation}
The time constant $\tau_{RC} = RC$ determines how quickly $\vC$ relaxes to
the steady-state value $\Vs$.

\subsubsection{First-Order RL Circuit}

Similarly, \eqref{eq:rl_ivp} yields:
\begin{equation}
  i(t) = \frac{\Vs}{R}\!\left(1 - e^{-(R/L)t}\right).
  \label{eq:rl_sol}
\end{equation}
The time constant is $\tau_{RL} = L/R$.

\subsubsection{Second-Order RLC Circuit: Characteristic Equation and Damping}

Seeking homogeneous solutions of the form $e^{rt}$ in \eqref{eq:rlc_ivp}
leads to the characteristic equation:
\begin{equation}
  LCr^2 + RCr + 1 = 0.
  \label{eq:char}
\end{equation}
Dividing by $LC$ and introducing the standard parameters
\begin{equation}
  \alpha = \frac{R}{2L}, \qquad \omega_0 = \frac{1}{\sqrt{LC}},
  \label{eq:params}
\end{equation}
the roots become $r = -\alpha \pm \sqrt{\alpha^2 - \omega_0^2}$.
The discriminant $\Delta = \alpha^2 - \omega_0^2$ determines damping behavior:

\begin{itemize}
  \item \textbf{Overdamped} ($\Delta > 0$, i.e.\ $R > 2\sqrt{L/C}$):
        two distinct real negative roots.
        The circuit returns to equilibrium without oscillating, but slowly.
        This occurs in high-resistance configurations such as lossy filters.
  \item \textbf{Critically damped} ($\Delta = 0$, i.e.\ $R = 2\sqrt{L/C}$):
        one repeated real root $r = -\alpha$.
        Fastest non-oscillatory return to equilibrium; used in precision
        instruments that must settle quickly.
  \item \textbf{Underdamped} ($\Delta < 0$, i.e.\ $R < 2\sqrt{L/C}$):
        two complex conjugate roots $r = -\alpha \pm j\omega_d$, where
        $\omega_d = \sqrt{\omega_0^2 - \alpha^2}$ is the damped natural
        frequency.
        The circuit oscillates with exponentially decaying amplitude; typical
        in resonant and tuned circuits.
\end{itemize}

\subsubsection{Second-Order RLC Circuit: General Analytic Solution}

\paragraph{Homogeneous solutions.}
For each damping case, the general homogeneous solution is:

\noindent\textit{Overdamped} ($r_1, r_2$ real, distinct):
\begin{equation}
  v_{C,h}(t) = A_1 e^{r_1 t} + A_2 e^{r_2 t}.
  \label{eq:over}
\end{equation}

\noindent\textit{Critically damped} ($r = -\alpha$ repeated):
\begin{equation}
  v_{C,h}(t) = (A_1 + A_2 t)\,e^{-\alpha t}.
  \label{eq:crit}
\end{equation}

\noindent\textit{Underdamped} (complex roots):
\begin{equation}
  v_{C,h}(t) = e^{-\alpha t}\!\left(A_1\cos\omega_d t + A_2\sin\omega_d t\right).
  \label{eq:under}
\end{equation}

\paragraph{Non-homogeneous solution.}
For a constant forcing $V(t) = \Vs$, the particular solution is simply
$v_{C,p} = \Vs$.
For time-varying inputs, two systematic methods apply:
(i)~\emph{Variation of parameters}, which replaces the constants $A_1, A_2$
with functions $u_1(t), u_2(t)$ determined by a two-equation system derived
from the ODE; or
(ii)~\emph{Laplace transforms}, which convert the ODE to an algebraic equation
$\mathcal{L}\{\vC\}(s) = \mathcal{L}\{V(t)\}(s) \cdot H(s)$, where
$H(s) = 1/(LCs^2 + RCs + 1)$ is the circuit transfer function,
followed by a partial-fraction expansion and inverse transform.

\paragraph{Applying initial conditions (underdamped case).}
The complete solution is $\vC(t) = v_{C,h}(t) + \Vs$.
Applying $\vC(0) = V_0$ and $\dot{\vC}(0) = 0$:
\begin{align}
  \vC(0) &= A_1 + \Vs = V_0 \;\implies\; A_1 = V_0 - \Vs, \label{eq:ic1}\\
  \dot{\vC}(0) &= -\alpha A_1 + \omega_d A_2 = 0 \;\implies\;
                 A_2 = \frac{\alpha(V_0 - \Vs)}{\omega_d}. \label{eq:ic2}
\end{align}
Substituting \eqref{eq:ic1}--\eqref{eq:ic2} into \eqref{eq:under} and
adding $\Vs$ gives the complete analytic solution for the underdamped case.
